\subsection{Testing strategy}

\subsubsection*{Current automated testing}

\Code{npm test} executes \Code{tests/backend.test.ts} through \Code{tsx} and Node's test runner. The five tests cover:
\begin{enumerate}
  \item ready lifecycle, ordering, duplicate/retry behavior, counter edit, closure, and number reuse;
  \item malformed payloads and media/settings validation without unintended mutation;
  \item SQLite persistence of turns, ordering, idempotency, media, and settings after reopen;
  \item local catalog discovery and rejection of unlisted/traversal-like paths;
  \item a real Nest HTTP server with two WebSocket clients, state/announcement broadcast, recall ordering, media broadcast, reconnect, close, and health.
\end{enumerate}

\Code{node tooling/audio/verify-audio.mjs} is a standalone asset check, not part of \Code{npm test}. TypeScript compilation is enforced by \Code{npm run build} but no type-check-only npm script exists.

\subsubsection*{Current manual testing evidence}

The superseded redesign note claimed manual viewport, offline, media, YouTube fallback, queue action, and audio activation checks. Those results are historical and not reproducible from a committed test script/report. Actual YouTube playback was explicitly unverified. They are not treated as current automated assurance.

\subsubsection*{Critical untested workflows}

Public audio activation/playback, React rendering/actions, public empty state, pagination, media duck/resume, autoplay policy, multiple real screens, browser compatibility, Electron install/upgrade, LAN exposure, corrupted/locked database recovery, large queues/history, settings UI/server-boundary mismatch, accessibility, and backup/restore.

\subsubsection*{Recommended layers}

Retain service/integration tests; add component tests for validation and controls, browser tests for two synchronized pages and audio activation with mocked media, packaged-app smoke tests, and operational recovery tests. Add load/security tests only after requirements are defined.

\subsubsection*{How tests are run}

\begin{CodeBlock}
npm ci
npm test
npx tsc --noEmit -p tsconfig.json
npx tsc --noEmit -p tsconfig.server.json
node tooling/audio/verify-audio.mjs
npm run build
\end{CodeBlock}

The final build command writes generated output; the preceding TypeScript commands are non-emitting checks.
